\section{从全连接网络到卷积：空间结构的利用}

\subsection{全连接网络处理图像的问题}

在第一讲中，我们学习了全连接神经网络（Fully Connected Neural Network）。如果直接用全连接网络处理图像，必须先将二维图像\textbf{展平}（flatten）为一维向量：

\begin{figure}[H]
\centering
\includegraphics[width=0.9\textwidth]{slides-images/slide-22.jpg}
\caption{全连接网络处理图像的问题：展平操作破坏了所有空间信息，且参数量巨大。\protect\footnotemark}
\end{figure}
\footnotetext{Slides 第22页。}

\begin{warningbox}{全连接网络用于图像的两大缺陷}
\begin{enumerate}
    \item \textbf{空间信息完全丧失}：将二维图像展平为一维向量后，像素之间的空间关系被完全破坏。相邻像素在展平后可能相距甚远。
    \item \textbf{参数量爆炸}：每个像素都连接到隐藏层的每个神经元。例如，一张 $256 \times 256$ 的灰度图像有 65,536 个像素，如果隐藏层有 1,024 个神经元，仅第一层就需要 $65,536 \times 1,024 \approx 6700$ 万个参数。
\end{enumerate}
\end{warningbox}

关键问题是：\textbf{如何利用输入的空间结构来指导网络架构设计？}

\subsection{局部连接的思想}

解决方案的核心思想很简单：不再让每个神经元连接到整幅图像的所有像素，而是只连接到输入图像的一个\textbf{局部区域}（local patch）。

\begin{figure}[H]
\centering
\includegraphics[width=0.85\textwidth]{slides-images/slide-23.jpg}
\caption{利用空间结构：将输入图像的局部区域（patch）连接到隐藏层中的单个神经元。每个神经元只"看到"图像的一小块区域。\protect\footnotemark}
\end{figure}
\footnotetext{Slides 第23页。}

然后，使用\textbf{滑动窗口}（sliding window）的方式，让同一个权重矩阵在图像的不同位置上重复应用。这就是\textbf{卷积}（Convolution）操作的核心思想。

\begin{figure}[H]
\centering
\includegraphics[width=0.85\textwidth]{slides-images/slide-24.jpg}
\caption{滑动窗口：将输入的局部区域连接到隐藏层神经元，通过滑动窗口定义连接。关键问题：如何设置权重以检测特定特征？\protect\footnotemark}
\end{figure}
\footnotetext{Slides 第24页。}

\begin{importantbox}{卷积操作的三个核心特性}
\begin{enumerate}
    \item \textbf{局部特征提取}：使用一组权重（滤波器/filter）提取\textbf{局部特征}
    \item \textbf{多滤波器}：使用\textbf{多个滤波器}提取不同类型的特征
    \item \textbf{参数共享}：每个滤波器的参数在空间上\textbf{共享}（spatial sharing），即同一个滤波器应用于图像的所有位置
\end{enumerate}
\end{importantbox}

\subsection{本章小结}

全连接网络将图像展平会丧失空间信息且参数量爆炸。卷积的核心思想是局部连接加参数共享，既保留了空间结构，又大幅减少了参数数量。

\newpage
